%% ═══════════════════════════════════════════════════════════════════
%% §5 — Signatures of even unimodular forms.
%%
%% Drafted by a section agent that read the Lean substrate in full
%% (RokhlinHMRankFour, RokhlinHMDischarge, RokhlinFromHM, VanDerBlijReduction,
%% RokhlinClassification, ThetaDefiniteDischarge, ThetaModularWeight,
%% SpinRokhlinInterface, LatticeSignature, E8Signature, E8Literal, E8Lattice,
%% RokhlinArfNoGo, AlgebraicRokhlin). Extracted verbatim from that agent's
%% transcript, retained at papers/D2/_sections/signatures-section.md; it is
%% included by paper_draft.tex and is not to be rewritten without re-reading
%% the substrate.
%%
%% ⚠ The four `% TODO:` markers below are the honest record: van der Blij,
%% Milnor–Husemoller, Rokhlin 1952, Freedman 1982, Freedman–Kirby and the
%% Atiyah–Singer index parity are NOT held in full text in the project's
%% primary-source cache, so the attributions resting on them are unverified.
%% They are LaTeX comments and do not render. DO NOT delete them to make a
%% readiness gate go green; a bundle carrying TODOs is meant to stay amber.
%%
%% The section deliberately does NOT cite E8Lattice.e8_det_one /
%% e8_diagonal_all_two / e8_diagonal_even (proved by native_decide, hence not
%% kernel-pure, and with no in-tree theorem linking CartanMatrix.E₈ to the
%% literal matrix E8lit that the chain actually runs on), nor e8_sigma_div_8,
%% rokhlin_gap or char_sq_valid_e8. It cites the kernel-pure E8lit chain.
%% ═══════════════════════════════════════════════════════════════════

\section{Signatures of even unimodular forms}
\label{sec:signatures}

\subsection{The result, and what is new about the route}
\label{sec:sig-result}

Let $M$ be a symmetric $n\times n$ integer matrix that is
\emph{unimodular} ($\det M = \pm 1$) and \emph{even} (every diagonal
entry $M_{ii}$ is divisible by $2$, equivalently $b(x,x)\in 2\mathbb{Z}$
for every $x\in\mathbb{Z}^n$). Write $\sigma(M)$ for its signature, the
number of positive minus the number of negative eigenvalues of the real
form $M\otimes\mathbb{R}$. The main result of this section is the
following.

\begin{theorem}[$8\mid\sigma$, unconditional]
\label{thm:eight-dvd}
For every $n$ and every even unimodular $M \in
\mathrm{Mat}_{n\times n}(\mathbb{Z})$, $\;8 \mid \sigma(M)$.
\end{theorem}

The statement is classical. It is van der Blij's congruence for even
forms, and it is standard textbook material on symmetric bilinear
forms~\cite{vanderBlij1959,MilnorHusemoller1973}.
% TODO: neither vanderBlij1959 nor MilnorHusemoller1973 was available to
% the author in full text; the attribution of Theorem 1 to these sources,
% and the characterisation of the classical proof as proceeding through
% characteristic vectors, are unverified against the primary literature.
No referee needs to be told that $8$ divides the signature of an even
unimodular form. What is new is the \emph{route}, and the route is what
this section is about.

Three properties distinguish it. First, it is complete inside a proof
kernel. Theorem~\ref{thm:eight-dvd} is the Lean~4 declaration
\texttt{eight\_dvd\_latticeSig}, whose statement is exactly the display
above with $\sigma$ interpreted as the honest Sylvester signature of the
real form and not as a free integer parameter, and whose proof carries no
hypothesis, no project-local axiom and no appeal to native compilation.
Two classical inputs were once carried as explicit hypotheses of the
reduction, namely a local-global statement for indefinite forms and a
modularity statement for definite ones. Both are now theorems in the same
library, so the divisibility is unconditional rather than reduced.

Second, the route is non-circular in a sense that can be checked
mechanically rather than argued informally. It does not compute the
four-dimensional spin bordism group, so it consumes no structure theory
of the spin cobordism ring~\cite{ABP1967}; it consumes neither Rokhlin's
theorem nor any statement equivalent to it; and it consumes no
divisibility statement about $\sigma$ at any point. The reader can
confirm this from the dependency graph rather than from our word for it:
every declaration named in this section is reachable from the
even-unimodularity predicate together with real quadratic-form theory,
$p$-adic analysis and the theory of the lattice theta function, and from
nothing else. This matters because the divisibility statements in this
area are close enough to one another that a proof of one from another is
easy to write down and hard to notice.

Third, the argument is sharp at the boundary where sharpness is usually
asserted rather than shown. The bound of Theorem~\ref{thm:eight-dvd}
cannot be improved to $16$: the $E_8$ form is even unimodular with
$\sigma = 8$, and both facts are machine-checked
(\texttt{e8lit\_even\_unimodular} and \texttt{e8lit\_latticeSig}). Every
strengthening to $16 \mid \sigma$ in the literature is therefore a
statement about $4$-manifolds and not about forms, and any formalisation
of such a strengthening must expose the extra input.
Section~\ref{sec:sig-residual} isolates that input, and
Section~\ref{sec:sig-arf} closes off the most natural attempt to remove
it.

\paragraph{What this does and does not do for the rest of the paper.}
Theorem~\ref{thm:eight-dvd} does \emph{not} enter the derivation of
$N_f \in 3\mathbb{Z}$. That derivation runs entirely through the chiral
central charge and the modular framing constraint, and it would stand
unchanged if this section were deleted. Nor does
Theorem~\ref{thm:eight-dvd} enter the $\mathbb{Z}_{16}$ argument of
Section~\ref{sec:z16-drinfeld}. The result is a second, independent
contribution of the present work, unconditional where the generation
constraint is conditional on a proposal of another author. It is
presented here because it is the unconditional mathematics in this paper
and because the integer $8$ it produces is the same integer that governs
the signature of the intersection form of a spin $4$-manifold, which is
the object the anomaly discussion of Sections~\ref{sec:modular}
and~\ref{sec:z16-drinfeld} is ultimately about. The connection is one of
subject matter, not of logical dependence, and we do not claim more.

\subsection{The signature calculus, and the reduction}
\label{sec:sig-calculus}

\paragraph{The signature, grounded.}
A formalisation of a signature theorem is only as good as its signature.
We define
\begin{equation}
\sigma(M) \;:=\; \mathrm{sig}^{+}\!\left(M\otimes\mathbb{R}\right)
\;-\; \mathrm{sig}^{-}\!\left(M\otimes\mathbb{R}\right),
\label{eq:latticesig}
\end{equation}
where $\mathrm{sig}^{\pm}$ are the positive and negative inertia indices
of the real quadratic form $x \mapsto x^{\mathsf T} M x$ on
$\mathbb{R}^n$, taken from the general theory of Sylvester's law of
inertia rather than introduced for this purpose
(\texttt{LatticeSignature.lean}). Two consequences are used below and
both are proved from~\eqref{eq:latticesig} directly: orientation reversal
$\sigma(-M) = -\sigma(M)$, and the rank bound $|\sigma(M)| \leq n$. When
$M\otimes\mathbb{R}$ is positive definite, $\sigma(M) = n$; when it is
negative definite, $\sigma(M) = -n$.

\paragraph{The calculus.}
Three structural laws of $\sigma$ carry the argument, and each is a
separate machine-checked statement.

\emph{Congruence invariance.} For $P$ with $\det P$ a unit in
$\mathbb{Z}$,
\begin{equation}
\sigma\!\left(P^{\mathsf T} M P\right) \;=\; \sigma(M)
\label{eq:sylvester}
\end{equation}
(\texttt{latticeSig\_congr}). This is Sylvester's law transported from
$\mathbb{R}$ to $\mathbb{Z}$-congruence, and it is what allows a
signature computed in one basis to be quoted in another. A companion
statement gives the same invariance under relabelling of the index type,
which is needed because a block decomposition produces a sum type rather
than $\mathrm{Fin}\,n$.

\emph{Block additivity.} For nondegenerate blocks $A$ and $B$,
\begin{equation}
\sigma(A \oplus B) \;=\; \sigma(A) + \sigma(B),
\label{eq:blockadd}
\end{equation}
where $A\oplus B$ is the block-diagonal matrix
(\texttt{latticeSigOf\_fromBlocks}). Nondegeneracy is not decoration.
The inertia indices of a degenerate form do not add across a direct sum
in the way~\eqref{eq:blockadd} asserts, because the radical is
distributed between the summands; the hypothesis is discharged for each
block that occurs here from unimodularity.

\emph{Generator values.} $\sigma(E_8) = 8$, $\sigma(-E_8) = -8$ and
$\sigma(H) = 0$, where $H = \left(\begin{smallmatrix} 0 & 1 \\ 1 & 0
\end{smallmatrix}\right)$ is the hyperbolic plane. The last two follow
from the first by orientation reversal and by direct computation. The
first is the only analytically nontrivial evaluation in this section,
because a signature is a count of eigenvalue signs and therefore a
statement about $\mathbb{R}$. We obtain it from an integral certificate
rather than from an eigenvalue computation. There is an explicit integer
matrix $C$, whose columns are twice the $E_8$ simple roots, with
\begin{equation}
C^{\mathsf T} C \;=\; 4\,E_8 .
\label{eq:cholesky}
\end{equation}
Both sides of~\eqref{eq:cholesky} are products of small integers, so the
identity reduces in the kernel by evaluation (\texttt{c8\_eq}). Cast to
$\mathbb{R}$ it exhibits $4E_8$ as a Gram matrix; $C$ is invertible
because $\det(C)^2 = \det(4E_8) = 4^8\det(E_8) \neq 0$, which uses
unimodularity and evaluates no $8\times 8$ determinant; hence $4E_8$ and
then $E_8$ are positive definite, and $\sigma(E_8) = 8$ follows from the
definite case of~\eqref{eq:latticesig}. The doubling in~\eqref{eq:cholesky}
is what makes the certificate integral: the $E_8$ root lattice has a
half-integer basis, so the naive $LDL^{\mathsf T}$ pivots are rational and
a rational certificate does not reduce.

Unimodularity and evenness of $E_8$ are established on the same explicit
matrix by exhibiting the integer inverse and checking
$E_8 \cdot E_8^{-1} = \mathbf{1}$, which is a few hundred small integer
multiplications, rather than by expanding the determinant
(\texttt{e8lit\_mul\_inv}, \texttt{e8lit\_even\_unimodular}).

\paragraph{The reduction.}
The classical route to Theorem~\ref{thm:eight-dvd} is the classification
of even unimodular lattices: every such $L$ is isomorphic to
\begin{equation}
E_8^{\,\oplus a} \;\oplus\; (-E_8)^{\,\oplus b} \;\oplus\; H^{\,\oplus c},
\qquad \sigma(L) = 8(a-b),
\label{eq:classification}
\end{equation}
after which $8\mid\sigma$ is immediate. The calculus above is exactly
what is needed to make~\eqref{eq:classification} yield the conclusion in
a kernel: each generator satisfies $8\mid\sigma$,
\eqref{eq:blockadd} propagates the divisibility across the direct sum, and
\eqref{eq:sylvester} transports it back along the isomorphism. Those
closure facts are shipped (\texttt{RokhlinClassification.lean}).

The proof we actually give does not use the normal
form~\eqref{eq:classification}. It uses the recursion that produces it,
which is weaker and therefore cheaper, and it is stated as follows.

\begin{proposition}[Reduction to two classical inputs]
\label{prop:reduction}
Suppose
\begin{itemize}
\item[{\rm[HM]}] every indefinite even unimodular integer form admits a
primitive isotropic vector, and
\item[{\rm[$\Theta$]}] every definite even unimodular integer form
satisfies $8\mid\sigma$.
\end{itemize}
Then $8\mid\sigma(M)$ for every even unimodular integer form $M$.
\end{proposition}

The proof is strong induction on the rank
(\texttt{eight\_dvd\_latticeSig\_of\_HM\_of\_Theta}). If $M$ is definite,
meaning $\mathrm{sig}^{+} = 0$ or $\mathrm{sig}^{-} = 0$, apply
[$\Theta$]. Otherwise $M$ is indefinite and [HM] supplies a primitive
$v$ with $v^{\mathsf T} M v = 0$. Primitivity and unimodularity give a
$w$ with $v^{\mathsf T} M w = 1$, and evenness upgrades $w$ to a genuine
hyperbolic partner: writing $w^{\mathsf T} M w = 2k$, which is possible
precisely because the form is even, the vector $w' = w - k v$ satisfies
\begin{equation}
v^{\mathsf T} M v = 0, \qquad v^{\mathsf T} M w' = 1, \qquad
w'^{\mathsf T} M w' = 0
\label{eq:hyperbolic-pair}
\end{equation}
(\texttt{exists\_hyperbolic\_pair}). The plane spanned by $v$ and $w'$ is
a copy of $H$ and splits off orthogonally. Passing to the Gram matrix
$M'$ of the orthogonal complement, one has
\begin{equation}
\sigma(M) = \sigma(M')
\label{eq:split}
\end{equation}
because $\sigma(H) = 0$ (\texttt{latticeSig\_split}), and $M'$ is again
even unimodular, of rank $n-2$
(\texttt{residGram\_evenUnimodular}). The induction hypothesis applies.

The recursion terminates at the definite forms and at rank $0$, and the
small ranks it passes through are not vacuous, so [HM] is needed at every
rank the recursion can reach. This is the point at which a reduction
becomes a proof, and Sections~\ref{sec:sig-indefinite}
and~\ref{sec:sig-definite} discharge the two inputs.

\paragraph{Scope of the classification claim.}
We are claiming Proposition~\ref{prop:reduction} together with the
discharge of [HM] and [$\Theta$], not the normal
form~\eqref{eq:classification} itself. The existence statement in
\eqref{eq:classification} is not in our library. Its
signature-zero case is (a form all of whose blocks are hyperbolic
planes), but the identification of the definite part with copies of
$E_8$ is not; the definite branch below stops at the divisibility
$8 \mid \mathrm{rank}$, which is all that $8\mid\sigma$ requires. A
reader who wants~\eqref{eq:classification} will not find it here, and we
do not use it.

\subsection{The indefinite branch, and rank four}
\label{sec:sig-indefinite}

\begin{proposition}[{[HM]} discharged at every rank]
\label{prop:hm}
Every indefinite even unimodular integer form admits a nonzero integer
isotropic vector, and hence a primitive one.
\end{proposition}

The upgrade from nonzero to primitive is content extraction: divide by
the greatest common divisor of the coordinates, which does not disturb
isotropy (\texttt{exists\_primitive\_isotropic\_of\_isotropic}). The
substance is the nonzero statement (\texttt{hasWeakIsotropicVector}), and
it splits by rank.

\emph{Ranks one and three do not occur.} A $1\times 1$ even unimodular
form would have $M_{00}$ simultaneously even and equal to $\pm 1$. For a
$3\times 3$ symmetric matrix with even diagonal, every term of the
Leibniz expansion is even, so $2 \mid \det M$, contradicting
unimodularity (\texttt{not\_evenUnimodular\_one},
\texttt{not\_evenUnimodular\_three}). These are the rank-$1$ and rank-$3$
instances of ``even unimodular implies even rank'', proved directly so as
to avoid characteristic-two Pfaffian machinery, and they are elementary
in the strict sense that they use no divisibility property of $\sigma$.

\emph{Rank two is explicit.} With both diagonal entries even,
$\det M = M_{00}M_{11} - M_{01}^2 \equiv -M_{01}^2 \pmod 4$, which lies
in $\{0,3\}$, so $\det M = 1$ is impossible and $\det M = -1$
(\texttt{det\_eq\_neg\_one\_of\_evenUnimodular\_two}). The vector
$(1 - M_{01},\, M_{00})$, or $(1,0)$ when $M_{00} = 0$, is then isotropic
by direct computation, since its value is $M_{00}(1 + \det M)$. Note that
no real-signature input is used: every even unimodular binary form is
indefinite.

\emph{Rank at least five is generic.} At rank $\geq 5$ a diagonal form
over $\mathbb{Q}_p$ is isotropic for every $p$, with no hypothesis beyond
nonvanishing coefficients, and over $\mathbb{R}$ isotropy follows from
indefiniteness. The local-global spine then applies
(\texttt{diag\_indefinite\_rat\_zero},
\texttt{weakIsotropic\_of\_five\_le}). Even unimodularity is not required
at this rank. Two pieces of plumbing make this usable on a Gram matrix
rather than a diagonal form. The form is diagonalised over $\mathbb{Q}$
and the congruence is exhibited explicitly so that it can be cast to every
completion; and indefiniteness, which is stated over $\mathbb{R}$, is
transferred to $\mathbb{Q}$ by density, since a real vector on which the
form is positive lies in an open set and therefore contains a rational
point (\texttt{sigPos\_cast\_pos}, \texttt{sigNeg\_cast\_pos}).

\emph{Rank four is the frontier.} It is the only rank at which local
isotropy can genuinely fail: anisotropic quaternary forms exist over every
$\mathbb{Q}_p$, so nothing is automatic and the even unimodular structure
must be used. We give the argument in full.

\begin{proposition}[Rank-four weak isotropy]
\label{prop:rank4}
Every indefinite even unimodular $4\times 4$ integer form $A$ admits a
nonzero integer vector $v$ with $v^{\mathsf T} A v = 0$.
\end{proposition}

\emph{Step 1: the discriminant is a square.} We first show
$\det A = +1$, excluding the shapes with $\det A = -1$ and
$\sigma = \pm 2$. Write the diagonal entries as $2g_i$ and the
off-diagonal ones as $a = A_{01}$, $b = A_{02}$, $c = A_{03}$,
$d = A_{12}$, $e = A_{13}$, $f = A_{23}$. Reduce the Leibniz expansion
modulo $4$. Every monomial containing at least two diagonal factors is
divisible by $4$ outright. Every monomial containing exactly one diagonal
factor corresponds to a permutation with exactly one fixed point, hence to
a $3$-cycle; the $3$-cycle and its inverse contribute the same product
with the same sign, so those monomials pair off and their sum acquires the
missing second factor of $2$. What survives is
\begin{equation}
\det A \;\equiv\; \left(af - be + cd\right)^2 \pmod 4
\label{eq:det4}
\end{equation}
(\texttt{det4\_evenDiag\_sq}). A square modulo $4$ lies in $\{0,1\}$, and
$-1 \equiv 3$ does not, so unimodularity forces $\det A = 1$
(\texttt{det\_eq\_one\_of\_evenUnimodular\_four}).

Now diagonalise over $\mathbb{Q}$, writing $A = P^{\mathsf T}
\mathrm{diag}(w)\, P$ with $\det P$ a unit. Taking determinants,
$1 = \det(P)^2 \prod_i w_i$, so
\begin{equation}
\prod_{i} w_i \;=\; \left(\det P\right)^{-2}
\label{eq:sqdisc}
\end{equation}
is a square (\texttt{isSquare\_prod\_weights}). Clearing denominators
coordinatewise, the integral diagonal form $\sum_i d_i x_i^2$ with
$d_i = \mathrm{num}(w_i)\cdot\mathrm{den}(w_i)$ also has square
discriminant. This is the structural payoff of Step~1, and it is what
makes the rest of the argument possible: a quaternary form of square
discriminant is controlled by a \emph{single} binary Hilbert symbol,
whereas a general quaternary form is not.

\emph{Step 2: the real place.} Indefiniteness gives a positive and a
negative coefficient among the $d_i$, and a two-variable combination of
the corresponding coordinates is isotropic over $\mathbb{R}$
(\texttt{diag\_real\_isotropic\_of\_signs}).

\emph{Step 3: odd primes, by Chevalley--Warning and Hensel.} Fix an odd
prime $p$. Since $\det A = 1$ is a $p$-adic unit, $A$ is
$\mathbb{Z}_p$-unimodular, and its reduction $\bar A$ modulo $p$ is a
rank-four form over $\mathbb{F}_p$ with $\det\bar A \neq 0$. A quadratic
form in at least three variables over a finite field has a nontrivial
zero, because the number of solutions of a polynomial system whose total
degree is smaller than the number of variables is divisible by the
characteristic, and the trivial solution is one of them; this is the
Chevalley--Warning theorem, and it is available for an arbitrary finite
field and any rank at least three
(\texttt{finite\_field\_form\_isotropic}). Let $\bar x \neq 0$ be such a
zero.

The zero is \emph{simple}. Its gradient is $\bar A\bar x$, and if that
vanished then $\bar x$ would lie in the kernel of $\bar A$, contradicting
$\det \bar A \neq 0$. Choose a coordinate $j$ with
$(\bar A \bar x)_j \neq 0$, lift $\bar x$ arbitrarily to $x_0$ over
$\mathbb{Z}_p$, and vary only coordinate $j$. Setting $y$ equal to $x_0$
with its $j$-th coordinate zeroed, the form becomes a univariate quadratic
in that coordinate:
\begin{equation}
\left(y + t\,e_j\right)^{\mathsf T}\! A \left(y + t\,e_j\right)
= A_{jj}\,t^2 + 2\,(Ay)_j\, t + y^{\mathsf T} A y ,
\label{eq:univariate}
\end{equation}
the cross terms coinciding by symmetry of $A$
(\texttt{mulVec\_update\_coord\_quadratic}). Call this polynomial $F$.
At $t = a := (x_0)_j$ one has $F(a) = x_0^{\mathsf T} A x_0$, which is
divisible by $p$ by construction, and
$F'(a) = 2\,(A x_0)_j$, which is a $p$-adic unit because $p$ is odd and
$(\bar A\bar x)_j \neq 0$. Hence $|F(a)|_p < |F'(a)|_p^2 = 1$, and
Hensel's lemma produces an exact zero $z \in \mathbb{Z}_p$ of $F$. The
vector $y + z\,e_j$ is a nonzero $\mathbb{Z}_p$-isotropic vector for $A$,
and therefore a $\mathbb{Q}_p$-isotropic one
(\texttt{isotropic\_padicInt\_of\_unit\_det}). Transporting through the
explicit congruence of Step~1 gives isotropy of the diagonal form at every
odd $p$.

\emph{Step 4: the prime two, by reciprocity.} No local computation at
$p = 2$ is performed. Instead, the square-discriminant hypothesis
converts local isotropy at a place $v$ into the vanishing of a single
binary Hilbert symbol. Concretely, if $abcd = s^2$ with all entries
nonzero, then over any field
\begin{equation}
a x^2 + b y^2 - c z^2 - d w^2 = 0
\;\;\Longleftrightarrow\;\;
z^2 = -ab\,x^2 + ac\,y^2
\label{eq:quaternary-ternary}
\end{equation}
has a nontrivial solution (\texttt{quaternary\_sqdisc\_iso\_iff\_ternary}),
the right-hand side being the norm form of the quaternion algebra attached
to the pair $(\alpha,\beta) = (-ab,\,ac)$. Isotropy at $v$ is thus
equivalent to $(\alpha,\beta)_v = 1$. Steps~2 and~3 give
$(\alpha,\beta)_\infty = 1$ and $(\alpha,\beta)_p = 1$ for every odd
prime. Hilbert's product formula
\begin{equation}
\prod_{v} \left(\alpha,\beta\right)_v \;=\; 1 ,
\label{eq:hilbert-product}
\end{equation}
the product running over the real place and all primes and being finitely
supported, then forces the remaining factor to be $1$
(\texttt{hilbertPrime\_two\_eq\_one\_of\_real\_odd}). The place $2$, which
is the hard place for quadratic forms, is obtained for free from the
others. We emphasise that~\eqref{eq:hilbert-product} is itself a theorem
in our library and not an assumption: it is assembled from the
bimultiplicativity of the local symbols, quadratic reciprocity and the
supplementary laws (\texttt{hilbertGlobalProd\_eq\_one}).

\emph{Step 5: local to global.} With every place discharged, the
square-discriminant quaternary Hasse--Minkowski statement applies. It
routes through the ternary form on the right of
\eqref{eq:quaternary-ternary} and requires neither weak approximation nor
Dirichlet's theorem on primes in arithmetic progressions
(\texttt{quaternary\_sqdisc\_solvable\_of\_local\_no\_two},
\texttt{diag\_four\_solvable\_of\_local\_no\_two}). The resulting
nontrivial rational zero is cleared of denominators to an integral one
(\texttt{exists\_int\_isotropic\_of\_rat}), and transported back through
the congruence of Step~1 to $A$ itself. This proves
Proposition~\ref{prop:rank4} (\texttt{weakIsotropic\_rank\_four}).

Two remarks on this argument. First, it is the even unimodular hypothesis,
and nothing else, that supplies each local input: unimodularity gives the
$\mathbb{Z}_p$-unit determinant that Step~3 needs, and evenness gives
the mod-$4$ congruence~\eqref{eq:det4} that Step~1 needs. Second, the
argument is confined to rank four by design. At rank at least five the
local inputs are automatic and the even unimodular structure is not used
at all; at rank four it is used twice, and both uses are essential.

\subsection{The definite branch, by theta modularity}
\label{sec:sig-definite}

For a definite form, $\sigma = \pm\,\mathrm{rank}$, so $8\mid\sigma$ is
equivalent to $8\mid\mathrm{rank}$. The bridge from the inertia
hypothesis to definiteness is short. A unimodular form is nondegenerate,
so its radical is trivial; with $\mathrm{sig}^{-} = 0$, the identity
$\mathrm{sig}^{+} + \mathrm{sig}^{-} + \dim(\mathrm{radical}) =
\mathrm{rank}$ forces $\mathrm{sig}^{+}$ to be the full dimension, and a
form whose maximal positive subspace is everything is positive definite.
The case $\mathrm{sig}^{+} = 0$ is the mirror image, handled by negating
the form (\texttt{eight\_dvd\_latticeSig\_of\_definite}). So the content
is the following.

\begin{proposition}[Modular weight]
\label{prop:theta}
A positive-definite even unimodular integer form of rank $d$ has
$8 \mid d$.
\end{proposition}

The proof is the classical theta argument, carried out analytically. For
a real positive-definite $G$ set
\begin{equation}
\Theta_G(\tau) \;=\; \sum_{v \in \mathbb{Z}^d}
e^{\,\pi i \tau\, v^{\mathsf T} G v},
\qquad \mathrm{Im}\,\tau > 0 ,
\label{eq:theta}
\end{equation}
which converges absolutely (\texttt{latticeTheta}).

\emph{$T$-invariance.} If $G$ comes from an even integer form then
$v^{\mathsf T} G v \in 2\mathbb{Z}$ for every $v$, so each summand
in~\eqref{eq:theta} is unchanged under $\tau \mapsto \tau + 1$ and
\begin{equation}
\Theta_G(\tau + 1) = \Theta_G(\tau)
\label{eq:theta-T}
\end{equation}
termwise, with no summability argument required
(\texttt{latticeTheta\_T\_int}).

\emph{$S$-transformation.} Multivariate Poisson summation applied to the
Gaussian $v \mapsto e^{\pi i \tau v^{\mathsf T} G v}$ gives the
functional equation relating $\Theta_G$ to the theta function of the dual
form $G^{-1}$ (\texttt{latticeTheta\_S}). Two facts collapse it. A
positive-definite unimodular form has $\det = +1$, which removes the
$(\det G)^{-1/2}$ prefactor; and for an integral $A$ with unit
determinant, $A^{-1}$ is again integral and
$A^{-1} = (A^{-1})^{\mathsf T} A\, A^{-1}$, so the dual lattice is the
original lattice in a different unimodular basis and its theta function is
literally the same function (\texttt{latticeTheta\_inv\_eq}). What remains
is the self-transformation
\begin{equation}
\Theta_G\!\left(-1/\tau\right)
\;=\; \left(\tau/i\right)^{d/2}\,\Theta_G(\tau).
\label{eq:theta-S}
\end{equation}

\emph{The multiplier is forced.} Composing~\eqref{eq:theta-S}
with~\eqref{eq:theta-T} gives the combined step
$\Theta_G(-1/(w+1)) = ((w+1)/i)^{d/2}\,\Theta_G(w)$
(\texttt{theta\_ST}). Iterating it three times traverses the orbit
\begin{equation}
i \;\longmapsto\; \tfrac{i-1}{2} \;\longmapsto\; i-1
\;\longmapsto\; i ,
\label{eq:st-orbit}
\end{equation}
which closes because $(ST)^3$ is the identity in the modular group. The
accumulated multiplier must therefore be $1$. It can be divided out
because $\Theta_G(i)$ is a sum of positive reals and so is nonzero
(\texttt{latticeTheta\_I\_ne\_zero}), and the three factors combine to a
single power because each base has positive real part, keeping the
principal branch of the complex power additive
(\texttt{mul\_cpow\_of\_re\_pos}). The product of the three multipliers
along~\eqref{eq:st-orbit} is $(-i)^{d/2}$, so
\begin{equation}
(-i)^{d/2} = e^{-i\pi d/4} = 1
\quad\Longleftrightarrow\quad
8 \mid d
\label{eq:eight-dvd-rank}
\end{equation}
(\texttt{neg\_I\_cpow\_eq\_one\_iff\_eight\_dvd},
\texttt{eight\_dvd\_rank}). This proves Proposition~\ref{prop:theta}.

Proposition~\ref{prop:reduction} applied to Propositions~\ref{prop:hm}
and~\ref{prop:theta} proves Theorem~\ref{thm:eight-dvd}.

It is worth recording what the two branches cost. The definite branch is
analysis: Poisson summation, absolute convergence of a Gaussian lattice
sum, and the principal branch of $z^{s}$. The indefinite branch is
arithmetic: finite fields, $\mathbb{Z}_p$, Hensel lifting, and quadratic
reciprocity in the form of the Hilbert product formula. Neither branch
uses the other, and neither uses any topology.

\subsection{The residual factor of two, and why it is irreducible}
\label{sec:sig-residual}

The signature of the intersection form of a closed smooth spin
$4$-manifold is divisible by $16$, not merely by
$8$~\cite{Rokhlin1952}.
% TODO: Rokhlin1952 (Doklady Akad. Nauk SSSR 84, 221) was not available to
% the author; the statement of Rokhlin's theorem as quoted here is taken
% from its standard formulation and is unverified against the original.
The gap between Theorem~\ref{thm:eight-dvd} and that statement is exactly
a factor of two, and the purpose of this subsection is to say precisely
where the factor of two comes from and to prove that it cannot come from
the form.

We record the interface first. A closed smooth spin $4$-manifold
contributes two pieces of data to the argument and no others: its
intersection form on $H^2(M;\mathbb{Z})$, which is even by Wu's formula
and unimodular by Poincar\'e duality, and the divisibility
$2 \mid \sigma/8$. The Lean structure \texttt{SmoothSpinManifold4}
carries exactly these, with its signature \emph{defined} as
$\sigma(\text{form})$ in the sense of~\eqref{eq:latticesig} rather than
supplied as a free parameter. On that interface,
\begin{equation}
8 \mid \sigma(M)
\label{eq:interface-eight}
\end{equation}
is a derived theorem and not a carried field
(\texttt{SmoothSpinManifold4.eight\_dvd\_sig}); it is
Theorem~\ref{thm:eight-dvd} applied to the intersection form. Composing
it with the carried divisibility gives
\begin{equation}
8 \mid \sigma \;\wedge\; 2 \mid \sigma/8
\;\;\Longrightarrow\;\; 16 \mid \sigma
\label{eq:compose}
\end{equation}
(\texttt{SmoothSpinManifold4.rokhlin}). We state plainly that this second
statement is conditional on the interface field $2\mid\sigma/8$, and that
it must be: $16\mid\sigma$ is false for general even unimodular forms, as
$E_8$ witnesses. A formalisation that presented $16\mid\sigma$ as
unconditional at the level of forms would be proving something untrue.

\paragraph{What the residual factor is.}
The input $2\mid\sigma/8$ is the evenness of the $\hat A$-genus. For a
closed spin $4$-manifold the index of the Dirac operator equals
$-\sigma/8$, and in dimension four that index is even because the spinor
bundle carries a quaternionic structure commuting with the Dirac
operator, so its kernel and cokernel are quaternionic vector spaces.
Equivalently, and more geometrically, it is the vanishing of the Arf
invariant of the quadratic refinement carried by a smoothly embedded
characteristic surface, in the congruence of
Freedman--Kirby~\cite{FreedmanKirby1978},
\begin{equation}
\sigma(M) \;\equiv\; \Sigma\cdot\Sigma \;+\; 8\,\mathrm{Arf}(M,\Sigma)
\pmod{16},
\label{eq:freedman-kirby}
\end{equation}
in which $\Sigma \subset M$ is a smoothly embedded surface representing a
characteristic class and $\mathrm{Arf}(M,\Sigma)$ is the Arf invariant of
a quadratic refinement on $H_1(\Sigma;\mathbb{Z}/2)$. For a spin manifold
the class $0$ is characteristic, both correction terms drop out, and
\eqref{eq:freedman-kirby} reduces to $16\mid\sigma$.
% TODO: neither the Atiyah--Singer identification of the index with
% -sigma/8 and its quaternionic parity, nor the Freedman--Kirby
% congruence (5.14), could be checked against a primary source available
% to the author; both are stated from their standard formulations and are
% unverified.

\paragraph{Why it is irreducible.}
There are two statements here, of increasing strength, and we prove the
first and cite the second.

The algebraic statement is that no property of the intersection form can
supply the factor of two, and it follows from
Theorem~\ref{thm:eight-dvd} being sharp. The implication ``$M$ even
unimodular $\Rightarrow 16 \mid \sigma(M)$'' is false, refuted by $E_8$,
which is even unimodular with $\sigma = 8$; both halves of that
refutation are machine-checked. Consequently any correct derivation of
$16\mid\sigma$ must consume an input that is not a function of the
isomorphism class of the form. The interface above is honest precisely
because it exposes such an input rather than deriving it, and the previous
subsections show that everything else on the path is derivable.

The stronger statement is that the input cannot be supplied by the
topology of the underlying manifold either, but requires the smooth
structure. Freedman's classification of simply-connected closed
topological $4$-manifolds produces a closed simply-connected
\emph{topological} $4$-manifold whose intersection form is
$E_8$~\cite{Freedman1982}.
% TODO: Freedman (1982) was not available to the author; the existence and
% properties of the topological E8 manifold are stated from their standard
% formulation and are unverified against the primary source.
Its form is even, so it is spin in the sense available to a topological
manifold, and its signature is $8$. It therefore violates the conclusion
$16\mid\sigma$ while satisfying every hypothesis that does not mention
smoothness. Smoothness is what the Dirac operator needs in order to exist,
and it is what a smooth embedding of the characteristic surface
in~\eqref{eq:freedman-kirby} needs in order to be available. This is not
a defect of our formalisation to be repaired later. It is the theorem's
actual content, and the interface is built so that the content is visible.

\subsection{The lattice Arf invariant cannot be the residual factor}
\label{sec:sig-arf}

The most natural attempt to remove the residual input is to look for it in
the mod-two structure of the lattice itself. An even lattice $L$ carries a
canonical quadratic refinement of its mod-two intersection pairing,
\begin{equation}
\bar q : L/2L \longrightarrow \mathbb{Z}/2,
\qquad
\bar q(x) \;=\; \tfrac{1}{2}\,b(x,x) \bmod 2 ,
\label{eq:redquad}
\end{equation}
which is well defined precisely because $L$ is even
(\texttt{redQuad}). Since $\bar q$ is a quadratic refinement of a
nondegenerate symplectic form over $\mathbb{F}_2$, it has an Arf
invariant, readable off the sign of its Gauss sum through the genus-$g$
identity
\begin{equation}
\sum_{x} (-1)^{\bar q(x)} \;=\; 2^{\,g}\,(-1)^{\mathrm{Arf}(\bar q)} ,
\label{eq:gauss-arf}
\end{equation}
which is itself machine-checked (\texttt{gaussSum\_genus\_g}). Write
$\mathrm{Arf}(M)$ for the invariant so defined
(\texttt{arfOfForm}). The candidate bridge is
\begin{equation}
\sigma(M)/8 \;\equiv\; \mathrm{Arf}(M) \pmod 2
\qquad\text{for every even unimodular } M ,
\label{eq:arf-bridge}
\end{equation}
which would replace the geometric input of
Section~\ref{sec:sig-residual} by a finite computation on $L/2L$. It is
attractive: it has the right shape, it is finite, and it is the lattice
shadow of the genuinely correct geometric congruence
\eqref{eq:freedman-kirby}. It is also false.

\begin{theorem}[The lattice Arf bridge is refuted]
\label{thm:arf-nogo}
Statement~\eqref{eq:arf-bridge} is false. Explicitly, $E_8$ is even
unimodular with $\sigma(E_8)/8 = 1$ and $\mathrm{Arf}(E_8) = 0$.
\end{theorem}

The witness is computed rather than argued
(\texttt{lattice\_arf\_bridge\_refuted}). On $E_8/2E_8$, which has
$2^8 = 256$ elements,
\begin{equation}
\sum_{x \in E_8/2E_8} (-1)^{\bar q(x)} \;=\; +16 \;=\; +2^4 ,
\label{eq:e8-gauss}
\end{equation}
evaluated by kernel reduction and not by native compilation
(\texttt{gaussSum\_redQuad\_e8lit\_eq}). By~\eqref{eq:gauss-arf} with
$g=4$ the positive sign is the $\mathrm{Arf} = 0$ branch, so
$\mathrm{Arf}(E_8) = 0$, while $\sigma(E_8)/8 = 1$ is odd. The Gauss sum
also decomposes as $136 - 120$, and the count is recognisable: the $120$
classes on which $\bar q = 1$ are the $240$ roots of $E_8$ taken in
$\pm$ pairs. Thus~\eqref{eq:arf-bridge} would read $1 \equiv 0$.

Theorem~\ref{thm:arf-nogo} is not an isolated numerical accident, and
three structural facts show that the bridge cannot be repaired inside the
lattice category.

\emph{The lattice Arf invariant is constant.} Each generator of the even
unimodular family has vanishing Arf invariant
(\texttt{arfOfForm\_e8lit\_eq\_zero},
\texttt{arfOfForm\_neg\_e8lit\_eq\_zero},
\texttt{arfOfForm\_hyp\_eq\_zero}), and the Gauss sum is multiplicative
under orthogonal direct sums with every factor on the positive branch, so
$\mathrm{Arf}$ is additive and vanishes identically on even unimodular
forms. A constant function cannot equal $\sigma/8 \bmod 2$, which is not
constant.

\emph{It is orientation-blind.} Every coefficient of $\bar q$ is fixed
modulo two under negation of the form, because $-a \equiv a$ and
$-(2k)/2 = -k \equiv k$, so $\bar q(-M) = \bar q(M)$ and hence
$\mathrm{Arf}(-M) = \mathrm{Arf}(M)$
(\texttt{redQuad\_neg\_eq}, \texttt{arfOfForm\_neg\_eq}). But
$\sigma(-M) = -\sigma(M)$. No invariant derived from $\bar q$ can see the
sign of the signature, and $E_8$ against $-E_8$ makes that concrete: the
Arf invariants agree while the signatures are $8$ and $-8$.

\emph{The discriminant form is trivial.} For a unimodular lattice the
discriminant group $L^{*}/L$ vanishes. The entire finite quadratic-form
apparatus that produces signature congruences from Gauss sums therefore
has nothing to act on beyond the mod-two reduction, and it can certify
only $\sigma \equiv 0 \pmod 8$, which is Theorem~\ref{thm:eight-dvd} and
which we already have by other means. It is blind to $\sigma$ modulo
$16$. The refinement of $\bar q$ to values in $\mathbb{Z}/4$ does not
help, since on an even lattice $b(x,x)$ takes only even values and the
refinement collapses back to the same $\pm 1$ sum.

The distinction that~\eqref{eq:arf-bridge} elides is between two
different objects that share a name. The Arf invariant appearing in the
correct congruence~\eqref{eq:freedman-kirby} lives on
$H_1(\Sigma;\mathbb{Z}/2)$ for a smoothly embedded surface $\Sigma$ in a
$4$-manifold; the one appearing in~\eqref{eq:arf-bridge} lives on
$L/2L$ for the intersection lattice. The first depends on a smooth
embedding and the second does not, and the data determining the first is
not present in, and not recoverable from, the lattice.
Theorem~\ref{thm:arf-nogo} is the precise statement that the substitution
is unavailable, and it explains why the interface of
Section~\ref{sec:sig-residual} must keep $2\mid\sigma/8$ as an input.

The Gauss-sum machinery itself is correct and reusable: it proves
\eqref{eq:gauss-arf} and it computes $\mathrm{Arf}$ on any even form
presented in symplectic normal form. What it does not do, and provably
cannot do, is certify the factor of two.
